\section{All size conditions and the initial arithmetic size datum}
\label{sec:size-conditions}

\status{SZ1--SZ8 prove the complete literal \texttt{SizesAt}
package, with the nested uniform thresholds of \texttt{SzRow}.
Together with PN and RF they construct every field of \texttt{NumAt}
in SZ9. SZ10--SZ12 independently construct the initial
\texttt{SizeData\textsubscript{2}} from the stated arithmetic and
Harish--Chandra data, and transport it through the actual
\texttt{Over} construction with its exact host and word-code identities.
They neither choose an arithmetic host nor construct
\texttt{ChOK}, \texttt{Over}, or \texttt{RecOK}.
The initial zero numerical fields are subsequently replaced by RF's
canonical refill.}

\paragraph{Notation and the order of the choices.}
In SZ1--SZ9 fix $0<\epsilon<1/52$, a standard constant record $C$,
and the same actual $D,L$ as in PN and RF, with
$\operatorname{Val}(C,c_D,c_h,L)$. To avoid confusing the size record
with its logarithmic scale, write
\[
 \begin{gathered}
 \mathcal L=x^\epsilon,\quad \delta=x^{-\epsilon/8},\quad
 \theta=\epsilon/100,\quad z=\epsilon/J,\quad
 N=N(m),\quad r=k+m-1,\quad f=2n,\\
 \Delta=\lceil x^{c_D}\rceil_+,\quad
 h=2\lceil x^{c_h}\rceil_+,\quad S=4k\Delta,\quad \eta=mh,\\
 \Lambda=L.\mathrm{LaG},\quad d_2=|T_2|/n,\quad
 N_* =2n\delta,\quad B=n x^{-\epsilon},\\
 Q=Q_Sqd,\quad \ell_A=\frac{kx^\epsilon}{md_2},\quad
 D_*=nt^2(N+t)\mathcal L^2+kn^2(S+t+\mathcal L),\quad
 \ell_* =\frac{mb}{D_*}.
 \end{gathered}                                                   \tag{SZ.1}
\]
Here $b,a,b',w'',X',X''$ are exactly (RF.1). In particular,
\emph{$\ell_*$ has the numerator $mb$}, so that
$b/\ell_*=D_*/m$. All natural quantities in real inequalities are
coerced to $\mathbb R$.

Every assertion called uniform below has the order
\[
 \forall\epsilon\ \forall C\ \exists c_{D,0}\
 \forall c_D\geq c_{D,0}\ \exists c_{h,0}\
 \forall c_h\geq c_{h,0}\ \exists x_0\
 \forall D,L.                                                    \tag{SZ.2}
\]
The last quantifiers require the fixed $\epsilon$, $x\geq x_0$ and
\texttt{Val}. In particular $x_0$ cannot depend on $D,L$, the supplied
primes, $J$, or a chosen position. We take finite maxima of RF11's
power/exponential thresholds and those in PN. Since
$1\leq J\leq C_8:=C.P.Cpde$, every occurrence of $z$ in a threshold
can be bounded below by $\epsilon/C_8$. This keeps (SZ.2) intact.

\paragraph{SZ1: the seven basic rank, density and length conditions.}
Put $\alpha=1/2-\epsilon/16$, $\beta=2+8\epsilon$, and
$\gamma=9/2+10\epsilon$. The literal rank windows say
\[
 x^\beta\leq n\leq2x^\beta,\qquad
 x^\alpha\leq m\leq2x^\alpha,\qquad
 x^\gamma\leq k\leq2x^\gamma.                                  \tag{SZ.3}
\]
Thus $x^\theta\geq2$ proves \texttt{SizeN}, \texttt{SizeM}, and
\texttt{SizeK}; the square root in the middle window is exactly
$x^\alpha$. The rank axiom and $L.\mathrm{t\_win}$ give
$1\leq t\leq x^{C_t}\leq x^{C_t+\theta}$, proving
\texttt{SizeT}. The host's two bounds
$\delta n/2\leq|TT|\leq2\delta n$ give
\[
       x^{-\theta}\delta n\leq |TT|\leq N_*.
\]
This is \texttt{SizeTT}. Its other density clause says
$x^{-\epsilon}n\leq |T_2|$; divide by $n>0$ and use the actual
\texttt{typeDensity\_eq\_card\_div} to obtain
\texttt{SizeDe2}. Finally,
\[
 L_w=10\lceil405\log x\rceil_+
             +2\lceil30/\epsilon\rceil_++22
       \leq4050\log x+60/\epsilon+34.
\]
For example $\log x\geq5000+100/\epsilon$ makes the right side
at most $(\log x)^2$. This proves \texttt{SizeLw}, with no rounding
term discarded.

\paragraph{SZ2: all fourteen coupled rank conditions.}
We use $\log\mathrm{rd}_E\leq\mathcal L$, established independently
in SZ3 below. Set $\zeta=\epsilon/1000$. For large $x$, (SZ.3)
implies
\[
 m\leq x^{\alpha+\zeta},\quad f\leq x^{\beta+\zeta},\quad
 k\leq x^{\gamma+\zeta},\quad
 100+8|C_\Gamma|\leq x^\zeta.                                 \tag{SZ.4}
\]
The following table verifies the first eight clauses of
\texttt{SizeKRel}. A listed exponent bounds the whole left side
above by that power of $x$; the right exponent is a lower bound for
the right side. For the fourth row the displayed inequality is the
square of the source inequality. All numbers in the table are exact
terminating rational decimals.
\begin{center}\small
\begin{tabular}{p{0.47\textwidth}cc}
Required inequality & upper exponent & lower exponent\\\hline
$x^\theta m^5f\mathcal L x^\epsilon\leq k$
 & $4.5+9.7035\epsilon$ & $4.5+10\epsilon$\\
$8C_\Gamma m^5f\mathcal L\leq k$
 & $4.5+8.6945\epsilon$ & $4.5+10\epsilon$\\
$k\leq\sqrt{x}\,nB$
 & $4.5+10.001\epsilon$ & $4.5+15\epsilon$\\
$k\leq mB^2$
 & $4.5+10.001\epsilon$ & $4.5+13.9375\epsilon$\\
$m^2\mathcal L\leq k$
 & $1+0.877\epsilon$ & $4.5+10\epsilon$\\
$m\sqrt{x}\log x\leq k$
 & $2-0.0615\epsilon$ & $4.5+10\epsilon$\\
$x^\theta m^6\mathcal L\leq k\log k$
 & $3+0.641\epsilon$ & $4.5+10\epsilon$\\
$x^\theta m^3\mathcal L\log\mathrm{rd}_E
                  \leq k(\log k)^2$
 & $1.5+1.8255\epsilon$ & $4.5+10\epsilon$
\end{tabular}
\end{center}
The sixth row uses $\log x\leq x$, and the last two use
$\log k\geq1$. Each exponent gap is positive for the entire allowed
range of $\epsilon$. Since $B>0$, the fourth row is equivalent to
$\sqrt{k/m}\leq B$.

For the remaining six clauses, $m$ is even and
$r_S=m/2-1$. The elementary logarithmic comparison from RF11 gives
$\log f\leq x^{1/5}$, whereas (SZ.3) eventually gives
$m\geq4x^{2/5}$ and hence $r_S\geq x^{2/5}$.
Therefore
\[
 \log f\leq\sqrt{r_S},\qquad
 (\log f)^2\leq r_S\leq m^2\mathcal L.
\]
Next $\log\log b=\log\mathcal L\leq\mathcal L$ and
$x^\theta\mathcal L\leq f$, by comparing exponents in (SZ.3).
As $r\geq k\geq5$, $\log r,\log k\geq1$, proving
\[
 x^\theta\log\log b\leq rf\log r,
       \qquad x^\theta\mathcal L\leq f\log k.
\]
Lastly $f=2n\leq m^2n$ because $m\geq2$, and eventually
$f\leq k\leq r\leq r^2$ by (SZ.3). These are the literal last two
clauses. Thus no clause of \texttt{SizeKRel} is omitted.

\paragraph{SZ3: prime, conductor and root-discriminant conditions.}
Choose the $Y$ already supplied by \texttt{HostVal}; no prime is
chosen here. Its window gives $x^{1/J}\leq Y\leq x$, and
\texttt{PairRow} gives
\[
 \begin{split}
 \tfrac12 x^{z/5}&<\log q\leq x^{\epsilon/5},\\
 \tfrac12 x^{z/5}(\log x/J)^2
             &<\log d\leq x^{\epsilon/5}(\log x)^2,\\
 \log Q&\leq2x^{\epsilon/5}(\log x)^2.
 \end{split}                                                    \tag{SZ.5}
\]
At a uniform threshold $\log x\geq C_8$ and
$x^{\epsilon/(100C_8)}\geq2$. Thus $\log x/J\geq1$ and
$x^{z/100}\geq2$. The two lower bounds in (SZ.5) give
$x^{19z/100}\leq\log q,\log d$, and
$z/6<19z/100$ gives the additional $q$ lower bound.
Absorbing $(\log x)^2$ into $x^\theta$ gives both upper bounds
$x^{\epsilon/5+\theta}$.

The supplied host's actual derived discriminant facts, rather than a
new assertion about the existence of that host, give
\[
 1\leq\mathrm{rd}_E\leq\frac{Q_S}{2}\sqrt{qd}\leq Q.
                                                                    \tag{SZ.6}
\]
For completeness the lower host bound is
$2^{e+1-2^{-e}}\sqrt d$, which is at least one because $e\geq0$
and $d$ is a positive prime. The upper bound in (SZ.6) is exactly
$H.\mathrm{der.rd\_le}$. Consequently
\[
 0\leq\log\mathrm{rd}_E\leq2x^{\epsilon/5}(\log x)^2
                                     \leq\mathcal L.           \tag{SZ.7}
\]
Together with \texttt{Val.dlog}, these prove all eight clauses of
\texttt{SizeQD}. The host also gives
$p_{\mathrm{de}},p'_{\mathrm{de}}\leq Q^{C_8}$. Hence
\[
 \log p_{\mathrm{de}},\log p'_{\mathrm{de}}
      \leq2C_8x^{\epsilon/5}(\log x)^2\leq\mathcal L,
\]
which proves \texttt{SizePde}.

RF12 gives $t\leq x^{401}$. Take logarithms in the required
conductor inequality. The sufficient inequality is
\[
 36\log Q+150\log t
    \leq72x^{\epsilon/5}(\log x)^2+60150\log x
    \leq\tfrac12 x^\epsilon=\log\sqrt b.                       \tag{SZ.8}
\]
RF11 proves its last step with a uniform threshold. Exponentiation
proves $Q^{36}t^{150}\leq\sqrt b$, the complete
\texttt{SizeCond}. Here $Q$ is the literal conductor expression
in \texttt{szCond}; no further conductor identity is presumed.

\paragraph{SZ4: the exact lambda, rho and bottom windows.}
RF14 proves, from the rounded orbit count, constants $A_C>0$ and
uniform thresholds such that
\[
        x^\epsilon\leq\Lambda
                    \leq A_Cx^\epsilon(\log x)^2.
\]
Absorb $A_C(\log x)^2$ into $x^\theta$; this proves
\texttt{SizeLa}. The coefficient $x^\theta$ in
\texttt{SizeRho} is variable, so one must not substitute it into a
bound proved only for a fixed coefficient. Instead use (SZ.3) directly:
\[
 4N_*x^\epsilon\Lambda
       \leq16x^{2+9.885\epsilon},\qquad
 x^\theta\log x\Lambda
       \leq x^{1.02\epsilon}\log x.                            \tag{SZ.9}
\]
Each is at most $\frac12x^{2.5+10\epsilon}$ for all sufficiently
large $x$. Their sum is at most $k x^{-2}$ by (SZ.3).
Division by $k>0$ gives precisely \texttt{SizeRho}.
Finally \texttt{Val.bT\_le}, with $N_*=2n\delta$, is exactly
\texttt{SizeBot}; no enlargement of the bottom window is needed.

\paragraph{SZ5: the eta, node and strict-star conditions on the same datum.}
PN4 proves the five eta inequalities, PN5 the four coupled eta/cloud
inequalities, and PN14 the rank-bottom-cloud inequality. In the
literal order of \texttt{SizeEta}, the supplied entries are
\[
 (\mathrm{PN4}_3,\mathrm{PN4}_1,\mathrm{PN4}_2,
   \mathrm{PN4}_4,\mathrm{PN4}_5,
   \mathrm{PN5}_1,\mathrm{PN5}_2,\mathrm{PN5}_3,\mathrm{PN14}).
                                                                    \tag{SZ.10}
\]
Thus the ninth clause includes the exact $\ell_A$, $d_2$ and
$b'+S$ factors; it is not inferred just from a small-$\eta$ estimate.

PN12 proves the following strengthened bounds on $h=\eta/m$:
\[
 \begin{gathered}
 12D_*+2t+8\leq h,\quad kt\mathcal L\leq h,\quad
 A_r r^{19}f^5\mathcal L\leq h,\quad
 4N^3+2t+6\leq h,\\
 A_k k m^{9/2}\leq h,\quad
 2(r+mt)^3+t+3\leq h,\quad Nm^{10}\leq h,
 \qquad 0\leq A_r,A_k\leq x.
 \end{gathered}                                                    \tag{SZ.11}
\]
Take $A_r=A_k=x^\theta$, since $0<\theta<1$ and $x\geq1$.
Since $m\geq1$ and $D_*>0$,
$12b/\ell_*=12D_*/m\leq12D_*$. This proves all seven
clauses of \texttt{SizeNodes}, including the nonintegral power
$m^{9/2}$. Finally PN15 proves
\[
 b'+S+\eta+b/\ell_*<\ell_*,\qquad
                2\eta+t+8b/\ell_*+7<\ell_*.
\]
These are exactly \texttt{SizeLstar}. Each application keeps the
same $D,L,c_D,c_h$ and uses the maximum of the finitely many
thresholds; it is not a selection of different parameters for different
clauses.

\paragraph{SZ6: all eight spacing conditions.}
Take $c_D\geq412$ and $x\geq32$. RF12 gives
$m\leq x$, $k\leq x^5$, $N\leq x^{402}$; also
$\theta<1$ and $\Delta\geq x^{c_D}$. Therefore
\[
 Nm^{10}\leq x^{412}\leq\Delta,\qquad
 x^\theta mN\leq x^{404}\leq\Delta.
\]
The ceiling upper bound $\Delta\leq2x^{c_D}$ and RF11 imply
\[
        \log\Delta\leq\log2+c_D\log x
                                   \leq\tfrac12x^{\epsilon/4}.
                                                                    \tag{SZ.12}
\]
At the scale-comparison threshold for the fixed exponent $c_D+12$,
$x^{c_D+12}\leq w''\leq b'$. Hence
\[
             16xk^2\Delta
                    \leq32x^{c_D+11}\leq x^{c_D+12}\leq w''.
                                                                    \tag{SZ.13}
\]
As $m\leq k$, $x^\theta\leq x$, $S=4k\Delta$ and
$\Lambda\geq1$, (SZ.13) proves the four required inequalities
\[
 kS\leq w'',\quad x^\theta mk\Delta\leq w'',\quad
          4S\leq\Lambda b'/k,\quad
          4x^\theta k^2\Delta\leq\Lambda b'.
\]
For the last clause, the host's literal $v$ is
$v=\mathrm{rd}_E/2>0$. Thus $\log v\leq\log\mathrm{rd}_E
\leq\mathcal L\leq x$ by SZ3. Therefore
$x^\theta m\log v\leq x^3\leq k$. This remains true if
$\log v$ is negative; multiplication uses only positive factors.
All eight clauses of \texttt{SizeDel} are now proved.

\paragraph{SZ7: the actual weight envelope, prime scales and word count.}
For fixed $c_h$, RF11's exponential comparison gives, for
$c=\max(c_h,0)+10$,
\[
 \sqrt b\,x^c\leq w'',\qquad b'x^c\leq b,
       \qquad\eta+1\leq w''\leq b'.                           \tag{SZ.14}
\]
The last inequality follows also directly from
$\eta=2m\lceil x^{c_h}\rceil_+\leq4x^{1+\max(c_h,0)}$,
after increasing the threshold. The first part gives
\texttt{SizeW}. The proof of SZ4 shows that the bottom surplus
$2n\delta x^\epsilon\Lambda/k$ tends to zero; alternatively
$\Lambda\leq x^2$, $f=2n\leq x^3$ and $k\geq1$ bound it by
$x^6$. It follows from (SZ.14) that
\[
          b_T+b'/x\leq(x^6+2)b'\leq b.
\]
Use the three \emph{actual} size-record weight windows. At a top,
$b_T+b'/x\leq w_j\leq b_T+b'/x+\eta$; at a thick position,
$b\leq w_j\leq b+\eta+1$; at a thin position,
$w''\leq w_j\leq w''+\eta$. They give the common envelope
\[
               w''\leq w_j(\tau)\leq2b
             \quad(1\leq j\leq m,\ \hbox{every }\tau).       \tag{SZ.15}
\]
Since $X''w''=xb$ and $x\geq2$, (SZ.15) proves every
cross-position, cross-index clause of \texttt{SizeXpp}.

Write $\beta_\tau=(D.hc\,\tau).b_0$. The \texttt{HCVal}
formulas give $b_{m-1}=\beta_\tau+(m-1)h$ and
$\beta_\tau\geq w''$: top, thick and thin positions give,
respectively, $\beta_\tau\geq b'$, $\beta_\tau\geq b$, and
$\beta_\tau=\lceil w''\rceil$. Thus
\[
       b_{m-1}\leq\beta_\tau+\eta\leq2\beta_\tau,
\]
which is exactly \texttt{SizeBm}.

The host's interval $(kb)^8\leq g_i\leq2(kb)^8$ proves both
the lower bound and $g_i\leq2g_j$ in \texttt{SizeG}. Moreover
\[
 \log g_i\leq\log2+8\log k+8\mathcal L
          \leq40\log x+9\mathcal L+9
          \leq x^\theta\mathcal L.
\]
The last inequality is another fixed-logarithm comparison. RF16 gives
the complete \texttt{SizeWords} bound
$|\mathrm{Word}|\leq2^{k-3}$, and RF14 gives
$11\leq\mathrm{splitBound}$. These statements concern the supplied
primes and words; no new prime or code has been selected.

\paragraph{SZ8: assemble the complete literal \texttt{SizesAt}.}
The definition has twenty-three named size clauses and one final
splitting bound. The following exhaustive map records their suppliers:
\begin{center}\small
\begin{tabular}{p{0.68\textwidth}l}
Clauses & supplier\\\hline
$N,M,K,T,TT,De2,Lw$ & SZ1\\
$KRel$ & SZ2\\
$Cond,QD,Pde$ & SZ3\\
$La,Rho,Bot$ & SZ4\\
$Eta,Nodes,Lstar$ & SZ5\\
$Del$ & SZ6\\
$W,Bm,Xpp,G,Words$ and $11\leq\mathrm{splitBound}$ & SZ7
\end{tabular}
\end{center}
Take the maximum of the $c_D$ thresholds, then, for that $c_D$,
the maximum of the $c_h$ thresholds, and finally the maximum of the
$x$ thresholds. This is precisely the quantifier order (SZ.2).
All conjunctions refer to the same objects. The resulting statement is
the full \texttt{szRow\_sizesAt}. The separately defined
\texttt{SizeTu} does not occur in \texttt{SizesAt}; it has not
silently been added to this theorem.

\paragraph{SZ9: every field of \texttt{NumAt} and the density/code pair.}
Assume additionally the literal \texttt{HostArith}. PN16 constructs
an actual $P:\texttt{ParameterChoice'}\,D\,L$ with
\[
          \operatorname{parConsts}(P)=C.PK,
                       \qquad P.CnW=C.P.CnW=10.
\]
Apply RF15 to \emph{this same $P$}. It constructs an actual
$T:\texttt{ParamsT}\,D\,L\,P$ with its prescribed seven constants
and $T.eP=C.T.eP$. SZ8 supplies the complete \texttt{sizes} field.
The remaining fields of \texttt{NumAt} are supplied as follows:
\begin{center}\small
\begin{tabular}{p{0.44\textwidth}p{0.48\textwidth}}
Field or group of fields & proof\\\hline
$\rho_B\leq x^{-2}$; $x^\epsilon\leq\Lambda$;
$g_i\leq e^{9\mathcal L}$; both $p_{\mathrm{de}}<X'$ bounds
 & RF14, on the same $D,L$\\
$C_\Gamma,C_c,C_t,CT_2$ equalities; bottom upper bound;
exact log-five depth & the corresponding fields of \texttt{Val}\\
$ccK=\top$ & \texttt{HostVal}\\
$L_w=10\lceil(C.Ct+4)\log x\rceil_+
                 +2\lceil30/\epsilon\rceil_++22$
 & $C.Ct=401$ and the literal $L_w$ equality\\
All positions have the same $v$ column
 & $v_i=(k-2-i)4\Delta$, independent of the position\\
The three raw rank windows & \texttt{Val.ranks}, exactly (SZ.3)
\end{tabular}
\end{center}
Together with the six $P,T$ data and constant fields and the
\texttt{sizes} field, this accounts for all twenty-one fields of
the actual record. RF14 also gives \texttt{ChTT\_at} and
$\texttt{ChCode\_at}(L,3,C.P.CF)$. Taking the finite maxima once
more proves
\[
 \texttt{HostArith}(D.host)\Longrightarrow
 \operatorname{Nonempty}(\texttt{NumAt}(C,L))
 \ \land\ \texttt{ChTT\_at}(L)
 \ \land\ \texttt{ChCode\_at}(L,3,C.P.CF)                    \tag{SZ.16}
\]
as an actual \texttt{SzRow} theorem. Neither full
\texttt{ParameterChoice'} nor \texttt{SizesAt} remains an
unproved numerical premise of this implication. Existence of the
arithmetic/geometric $D$ is a different assertion.

\paragraph{SZ10: a torsion estimate and all thirteen initial-window clauses.}
We now prove the initial windows \emph{without assuming an existing
size datum or PN16}. Fix a scale $s$, actual $D$, top bottom-degrees
$b_T$, and the literal \texttt{RankVal}, \texttt{HostVal}, and
\texttt{HCVal}, with
\[
 c_D\geq0,\quad c_h\geq402,\qquad
                  b'\leq b_T\leq2b'.                         \tag{SZ.17}
\]
All thresholds here depend only on $\epsilon,c_D,c_h,C_8$.

First let $F$ be a number field of degree $d_F$. A root of unity
of exact order $a$ has cyclotomic degree $\varphi(a)\leq d_F$.
For $a=\prod p^{e_p}$,
\[
  \frac{\varphi(a)^2}{a}
          =\prod_{p\mid a}p^{e_p-2}(p-1)^2\geq\tfrac12.
                                                                    \tag{SZ.18}
\]
Each odd-prime factor is at least one; the factor at $2$, if present,
is at least $1/2$. Thus $a\leq2d_F^2$. This first proves that all
roots of unity form a finite group, since their orders are bounded.
A finite subgroup of a field's multiplicative group is cyclic, so its
generator shows that its cardinality is also at most $2d_F^2$.
For the supplied host $[E:\mathbb Q]=2f$, giving
\[
             |\mu(E)|\leq8f^2\leq x^7\leq h/2             \tag{SZ.19}
\]
after RF12 and $x\geq8$. This supplies the actual torsion bound used
by the thick-position interval.

RF11--RF12, without a size record, give at a common threshold
\[
 k+5\leq b'/x,\quad \eta\leq b,\quad b'\leq b,\quad
 6\leq b,\quad1\leq t\leq x^{401},\quad nN/m\leq x^{405}.
                                                                    \tag{SZ.20}
\]
Here the crude $f\leq x^3$, $m\leq x$, $N\leq x^{402}$
bounds are used only where needed. By \texttt{HCVal}, writing
$\beta_\tau$ for its integral base value,
\[
 b_j=\beta_\tau+jh,\quad e_j=\beta_\tau+jh+h/2,
       \qquad w_j=e_{j-1}\ (1\leq j\leq m).                  \tag{SZ.21}
\]
In particular
$\beta_\tau+1\leq w_j\leq\beta_\tau+\eta-h/2$.
We check all thirteen clauses of the source window theorem:
\begin{enumerate}
\item At a top,
$b_0-b_T=\lfloor b'/x\rfloor>(k-1)/2+3$ by (SZ.20).
\item For every position and $0\leq j<m$, $b_j\leq b^2$.
Indeed (SZ.17) and the three possible base values imply
$\beta_\tau\leq3b+2$: tops use $2b'+b'/x$, thick positions
use $\lceil b\rceil+|\mu(E)|$, and thin positions use
$\lceil w''\rceil$ with $w''\leq b'$. Equations
(SZ.19)--(SZ.21) give $b_j\leq4b+2\leq b^2$ for $b\geq6$.
\item At a thick position the integral base satisfies
$\lceil b\rceil\leq\beta_\tau<\lceil b\rceil+|\mu(E)|$.
Since both endpoints are integers and $|\mu(E)|\leq h/2$,
$b\leq\beta_\tau\leq b+h/2$. Consequently
$b\leq w_j\leq b+\eta+1$, the required thick window.
\item At a top,
$\beta_\tau=b_T+\lfloor b'/x\rfloor$. The bounds
$b'/x-1<\lfloor b'/x\rfloor\leq b'/x$ and $h/2\geq1$
give $b_T+b'/x\leq w_j\leq b_T+b'/x+\eta$.
\item At a thin position $\beta_\tau=\lceil w''\rceil$.
The same $h/2\geq1$ gives $w''\leq w_j\leq w''+\eta$.
\item Off the top set, $u_0-v_0=S+1/2$.
Since $S=4k\Delta\geq1$, it lies in $[S,2S]$.
\item For every $i<k-1$,
$v_i=(k-2-i)4\Delta$ lies in $[0,S]$.
\item For $i,j<m$,
$|b_i-e_j|=|(i-j)h-h/2|\leq\eta$.
\item The rank axiom gives $1\leq t$.
\item RF12 gives $t\leq x^{401}$.
\item Set $q_0=\lceil\log_5(nN/m)\rceil_+$. Since
$\log5\geq1$, (SZ.20) gives
$q_0\leq\lceil405\log x\rceil_+$. Thus
$10q_0+20\leq L_w$ by its exact formula.
\item $|T_2|>0$ follows from $|T_2|\geq x^{-\epsilon}n>0$.
\item $|T_2|\leq n$ follows from the host's actual
\texttt{second\_set\_range}: its integer labels lie in the range
of length $n_0\leq n$ (the condition $L_w\geq2$ follows from
the exact formula). No density assumption stronger than the source
one is substituted.
\end{enumerate}
These give the complete \texttt{windows\_of\_explicitParameters}
conclusion. The multiplicative-unit condition in \texttt{HCVal}
is retained for subsequent parameter construction; none of these
thirteen window estimates needs to manufacture it.

\paragraph{SZ11: construct the initial size record field by field.}
Assume the literal inputs to \texttt{exists\_sizeData\_explicit}:
\texttt{Std}, $c_D\geq0$, $c_h\geq402$, the arithmetic and
column data of SZ10, the bounded logarithmic derivative, the prescribed
$\Delta$, $CT_2$, word count and $\eta_F$, and
\[
 \log\Gamma_u=C_\Gamma m^5f\mathcal L,
 \qquad
 b'\leq b_T\leq b'+2n\delta x^\epsilon\Lambda b'/k,           \tag{SZ.22}
\]
where now $\Lambda=G.\mathrm{La}(C.GK,s)$. Also retain the source's
$a_T$ window, the two top $u$ endpoints, and the untwisted equality
$A_j=0$. These are actual constructor inputs, not conclusions of
this step.

There is no circular use of SZ4's existing-size hypothesis here.
RF14's independent orbit-count bound gives
$\Lambda\leq A_Cx^\epsilon(\log x)^2$ directly for the scale.
Using $\delta\leq1$ and the raw rank windows,
\[
 \frac{2n\delta x^\epsilon\Lambda}{k}
       \leq4A_C x^{-5/2}(\log x)^2\leq1.                    \tag{SZ.23}
\]
Thus (SZ.22) implies (SZ.17), and SZ10 applies. Increase the
threshold so that $C_8\leq x^\theta$. Define the initial record by
\[
 \begin{array}{c|l}
 \text{fields}&\text{chosen values}\\\hline
 \mathrm{sc},a_T,b_T&s,a_T,b_T\\
 \mathrm{supp},\mathrm{drop}&\varnothing,\varnothing\\
 \mathrm{gain},F,E,\tau,\mathrm{io}&0,0,0,0,0\\
 \mathrm{NuLat},\mathrm{lenMN}&\perp,0\\
 j_0&\lceil\log_5(nN/m)\rceil_+\\
 d_2,\Lambda,S,\eta&|T_2|/n,G.\mathrm{La}(C.GK,s),4k\Delta,mh\\
 C_c,C_t,C_\Gamma&C.Cc,C.Ct,C.CGam\\
 TT_w,\mathrm{wPart}&TT,1
 \end{array}                                                     \tag{SZ.24}
\]
Take the logarithmic-derivative field to be the actual
$\operatorname{dlogL}(D,1)=\operatorname{dlogH}(D.host)$.
The $a_T,b_T$ and top-$u$ fields are the supplied ones. SZ10 proves
the top-bottom separation, upper $b_j$ bound, three weight windows,
cloud gap, $v$ bounds, $b/e$ closeness, depth and $t$ bounds.
For each word's middle $u_i$, strict interlacing gives
$v_i<u_i<v_{i-1}$; since these $v$ columns are unchanged by the
word, $0\leq u_i\leq S$. This supplies \texttt{cloud\_u} on
its exact range $1\leq i$, $i+1<k$.

All support, conjugacy, non-inert and above-band implications for
the chosen zero fields are immediate equalities. The untwisted field
uses the supplied $A_j=0$ and $F=0$. The empty drop set excludes
every type prime, and $TT_w=TT$ proves its fullness condition.
The splitting tie is the host field. The other two ties follow from
\[
 o\leq C_8\leq x^\theta,\qquad
             Q_S\leq C_8n\leq4nx^\theta.
\]
SZ10 gives $0<|T_2|\leq n$, proving $0<d_2\leq1$ and its
defining density equality. The positive $S,\eta,C_\Gamma$ and
unit lower bound on $\mathrm{wPart}$ are immediate from the chosen
values and \texttt{Std}. Finally the exact
\texttt{Gam\_eq} is the supplied equality in (SZ.22).
This accounts for every data and proof field of
\texttt{SizeData\textsubscript{2}}, not merely its numerical subset.

The fifteen fields of \texttt{Val} now follow from the supplied rank,
host, column, derivative, spacing, $CT_2$, word and $\eta_F$ data,
together with the defining equalities (SZ.24) and the bottom window
(SZ.22). Thus there is an actual $L$ with
\[
      L.\mathrm{sc}=s,\quad L.a_T=a_T,\quad L.b_T=b_T,
                         \operatorname{Val}(C,c_D,c_h,L).       \tag{SZ.25}
\]
This is an initial record: its zero gains and lattices are legal
values for the constructor, and RF5 later replaces them by the
canonical refill. No final forcing or lattice estimate is asserted
from the zeros.

\paragraph{SZ12: transport through \texttt{Over} with exact labels and code.}
Supply the literal \texttt{ChOK} datum $P$, its level record $Lv$,
\[
 Lv.\Gamma=\Gamma_S(C,s,P.R),\qquad
             \texttt{Over}(D,P,Lv),\qquad\texttt{RecOK}(s,C,c_D,D).
                                                                    \tag{SZ.26}
\]
The existential decomposition of \texttt{Over} identifies $D$ with
its actual staged construction. Its host has the same rank, $TT$
labels, second set, $q,d,g,p_{\mathrm{de}},p'_{\mathrm{de}}$ and
distinguished translate as $P.H$; its level-prime set is the old
one united with the prescribed level primes. In particular the old
set is contained in the new one, rather than necessarily equal to it.
The \texttt{Pre.ofCh} construction proves that every member of this
union is in class \texttt{splitK}. The other three
\texttt{HostArith} clauses are exactly the arithmetic part of
\texttt{ChOK}. Hence \texttt{HostArith}(D.host) is supplied.

Use SZ11 with $a_T(t)=a(t.\mathrm{val})$ and
$b_T(t)=b(t.\mathrm{val})$. The rank, host, column, derivative,
$CT_2$, word, degree-window and top-$u$ inputs come from
\texttt{ChOK}; the spacing, $\eta_F$ and untwisted inputs come
from \texttt{RecOK}. The staged identity
$\Gamma_u=Lv.\Gamma$ and (SZ.26) give
$\log\Gamma_u=C_\Gamma m^5f\mathcal L$ by $\log(e^y)=y$.
Thus (SZ.25) yields $L.\mathrm{sc}=s$,
\texttt{Val}, and the exact \texttt{WtEq}: injectivity of the
frame-label map identifies any two top positions having the same
label, so their two degrees are the prescribed $a,b$ values.

For clarity the simultaneous identities proved by this construction
are
\[
 \begin{gathered}
 D.R=P.R,\quad
 D.host.TT.\mathrm{image}(\mathrm{val})
          =P.H.TT.\mathrm{image}(\mathrm{val}),\quad
 D.host.T_2=P.H.T_2,\\
 (q,d,g,p_{\mathrm{de}},p'_{\mathrm{de}})_D
          =(q,d,g,p_{\mathrm{de}},p'_{\mathrm{de}})_P,\quad
 P.H.\mathrm{levelPrimes}\subseteq D.host.\mathrm{levelPrimes},\\
 D.host.u_0.\mathrm{val}=P.H.u_0.\mathrm{val},\qquad
 (G.h_P(D,l)).\mathrm{val}=(h_{P.H}(l)).\mathrm{val}.
 \end{gathered}                                                    \tag{SZ.27}
\]
The word types are identified by an actual equivalence $e_W$, which
is the identity after unfolding the stages. For $l<k-2$, compare
the source's two expressions for the same $u_l$ entry: the
\texttt{u\_letter} formula and the \texttt{ChOK.code} formula.
Cancel their identical base terms to obtain
\[
 \Delta\,\mathbf1_{\mathrm{letter}(u,l)}
             =\Delta\,\mathbf1_{c(e_Wu,l)}.
\]
Since $\Delta>0$, the two Boolean bits agree. For $l\geq k-2$
both extended bits are false by definition. Hence, for every $u,l$,
\[
 \operatorname{wordBits}(D,u,l)=
 \begin{cases}c(e_Wu,l),&l<k-2,\\ \mathrm{false},&l\geq k-2.
 \end{cases}                                                      \tag{SZ.28}
\]
This proves the complete literal \texttt{exists\_sizeData\_over},
including the labels and code, and discharges the initial-size
supplier retained in RF10. The arithmetic column constructor
\texttt{chGood\_of\_host}, the level-prime constructor and the
geometric \texttt{RecExP}/\texttt{RecOK}/\texttt{LedOK} supply
remain distinct tasks. SZ9 can now be applied to this $L$ when its
uniform thresholds are included; RF's refill preserves
\texttt{Val}, so it can also be applied to that refilled datum.

\paragraph{SZ13: construct the actual ranks over the prescribed frame.}
The numerical theorems above assume an actual rank record. We can also
construct that record over any supplied frame $\phi$ whose order $n$
satisfies the required window
\[
             x^{2+8\epsilon}\leq n\leq2x^{2+8\epsilon}.
                                                                    \tag{SZ.29}
\]
This preserves the given group, labels, projection and distinguished
generator; no replacement of the frame is made. The uniform threshold
$x\geq16$ suffices. Put
\[
 A=\sqrt{x\delta}=x^{1/2-\epsilon/16},\qquad
 K=x^{9/2+10\epsilon},\qquad
          m=2\lceil A/2\rceil_+,\quad k=2\lceil K/2\rceil_+.
                                                                    \tag{SZ.30}
\]
For any $y\geq2$, $y\leq2\lceil y/2\rceil_+<y+2\leq2y$.
Here $2\leq A\leq x$, $x^4\leq K$, and
$x^2\leq n\leq2x^3$, using $0<\epsilon<1/52$.
Thus $k,m$ are even, their two rank windows hold, and
\[
       2\leq m,\quad m+3\leq k,\quad m\leq2n\leq k.
                                                                    \tag{SZ.31}
\]
For example $m\leq2x$, $2n\geq2x^2$, and
$2n\leq4x^3\leq x^4\leq k$ prove the last two comparisons;
$2x+3\leq x^4$ proves the remaining gap.

Set $f=2n$, $r=k+m-1$, and define
\[
 \begin{split}
 T&=\lceil x^{377}\rceil_+,\\
 W&=\left\lceil\frac{10r}{8}
       \left(f+\frac{f^2(k-1)^3\mathcal L^2}{m}\right)
                                                    \right\rceil_+,\\
 t&=(8f-1)r+4+8TW.
 \end{split}                                                    \tag{SZ.32}
\]
The last expression is a nonnegative integer and equals the source's
natural subtraction $8fr+4+8TW-r$, because $f\geq1$.
Since $r$ is odd and $8f-1$ is odd, $t$ is odd. Moreover
\[
             t\geq7r+4\geq r+5,\qquad 3m<t,
\]
where $r=k+m-1\geq m+4$. In particular $m\leq t+1$, so
$m^2\leq k+m(t+1)$. These are all the additional inequalities
in the literal \texttt{Ranks} record. Together with the unchanged
$\phi$ and (SZ.31), they therefore construct an actual $R$.

The definitions of $T$ and $8W$ are precisely $\mathrm{TOv}(s)$ and
$\mathrm{nWv}(R,s)$. There is no circular choice of $t$: the latter
expression depends on $n,f,k,m,r$ and the scale, all fixed before
$t$ is defined. Equation (SZ.32) gives
\[
             r+t=8fr+4+\mathrm{TOv}(s)\mathrm{nWv}(R,s).
                                                                    \tag{SZ.33}
\]
This, (SZ.29), and (SZ.30) are exactly
$\mathrm{RankVal}(s,R)$, with $R.\mathrm{toFrame}=\phi$.
Thus the literal \texttt{exists\_ranks\_rankVal\_frame} is
proved, uniformly over all such prescribed frames.

There is also an explicit cyclic-frame specialization. Given $B\geq1$,
choose the least integer $e\geq0$ with $2^e\geq B$. Minimality
gives $2^e\leq2B$. For $B=x^{2+8\epsilon}$, (SZ.29) therefore
holds with $n=2^e$, and $e\geq2$. The cyclic frame uses
$A=\mathbb Z/2^e\mathbb Z$, generator $1$, and the homomorphism
\[
 A\longrightarrow
  (\mathbb Z/2^{e+2}\mathbb Z)^\times/\langle-1\rangle,
                            \qquad j\longmapsto[5^j].       \tag{SZ.34}
\]
It is an isomorphism. Indeed
$v_2(5^{2^j}-1)=j+2$, by induction from
$(u^2-1)=(u-1)(u+1)$ and $u\equiv1\pmod4$.
Thus $5$ has order $2^e$ modulo $2^{e+2}$, and its powers never
equal $-1$ there because they are $1$ modulo $4$. The quotient
also has order $2^e$, so (SZ.34) is bijective.
Transport this additive group to the source's labels
$\mathrm{Fin}(2^e)$ by the residue-class bijection. Its projection
is surjective, and a multiple of its distinguished generator whose
projection is zero is itself zero, because the projection is
injective. These are exactly the frame axioms.
The chosen isomorphism is the source's \texttt{cycIso}: a homomorphism
from a cyclic group is determined by its value at $1$.
Consequently this is the literal \texttt{cyclicFrame}(e), and
(SZ.30)--(SZ.33) prove \texttt{exists\_ranks\_cyclic\_frame}.
This constructs rank and frame data only; the number field and
arithmetic host with those labels remain separate existence problems.
